% !TEX root = ../../../main.tex
% Block card: interrogation.cycle_shapes -- "Cycle shapes" (Part I, the blocks chapter).
\cardfield{What it does}
For a chain of repeating drops, takes each stretch from one drop to the next (a cycle),
stretches every cycle to the same length, and asks whether they all share one shape or fall
into a few kinds, and where in the chain the kind changes.

\cardfield{When to use it}
On a long, regular chain of drops (the ``sharkfin'' rhythms), to see whether the cycle's
shape stays the same over hours or switches. A chain is like a drummer whose tempo drifts:
stretching each beat to one length lets you compare the beats themselves. Put
\blockref{interrogation.event_shape} and \blockref{interrogation.trains} (with
\texttt{maxinterval} set to chain-scale gaps) before it and
\blockref{interrogation.intervals} after it.

\cardfield{What goes in, what comes out}
In go the events. They come out unchanged, the event that opens each cycle carrying that
cycle's measures: its length, its size, which shape group it is in, whether the shape changed
there, how long it sits on its floor before climbing, and where it is half way back. A
stacked picture of all the cycles comes with them.

\cardfield{Settings that matter}
\emph{max\_cycle\_s}: a cycle much longer than its neighbours is treated as a gap (a missed
drop), not a cycle. The automatic rule (0) leaves out cycles over 1.6 times the local typical
length; set a number if many real long cycles are being thrown away. \emph{min\_cycles}
(10) is the fewest usable cycles before any sorting happens. Bind the \emph{waveform} side
input to the root signal when the detector ran on a smoothed-out trace, so the slow climb of
each cycle is kept.

\cardfield{Watch out for}
More than one shape is claimed only when the split clearly beats the same split on stretches
of quiet background; with fewer than 20 usable cycles it can never find two kinds. A chain
that fills its whole span leaves no quiet background, and its verdicts are then marked
untested. The cycle size is in the recording's stored unit, not converted to millivolts.
